\documentclass[../../main.tex]{subfiles}
\graphicspath{{../../image/}{../../figure/}}
\begin{document}

\section{守恒律方程组的测度值解与补偿紧致性}\label{section:ch17-测度值解与补偿紧致}

考虑守恒律方程组
\begin{align}
\partial_t U+\partial_x F(U)=0,
\label{align:ch17-1731}
\end{align}
并设 $\{U_k\}$ 是 $\mathbb{R}^2$ 的开子集 $\Omega$ 中的一近似解序列，即
\begin{align}
\partial_t U_k+\partial_x F(U_k)\longrightarrow 0,\qquad k\to\infty,
\label{align:ch17-1732}
\end{align}
在 $\Omega$ 上的分布意义下成立。例如，$\{U_k\}$ 可由消失黏性方法导出，即 $U_k=U_{\mu_k}$，其中 $\mu_k\downarrow0$（$k\to\infty$），而 $U_\mu$ 是抛物方程组
\begin{align}
\partial_t U+\partial_x F(U)=\mu\,\partial_x^2U
\label{align:ch17-1733}
\end{align}
的解。

当 $\{U_k\}$ 含于 $L^\infty(\Omega;\mathbb{R}^n)$ 的一个有界集时，依照\cref{section:ch16-cauchy问题}的讨论，
可提取一个子列 $\{U_j\}$，它联系着一族 Young 概率测度 $\{\nu_{x,t}:(x,t)\in\Omega\}$，使得对任意连续函数 $h$，$h(U_j)\rightharpoonup\langle\nu,h\rangle$，当 $j\to\infty$ 时，在 $L^\infty$ 弱$^*$ 意义下成立。特别地，由\eqref{align:ch17-1732}，
\begin{align}
\partial_t\langle\nu_{x,t},U\rangle+\partial_x\langle\nu_{x,t},F(U)\rangle=0.
\label{align:ch17-1734}
\end{align}
于是可把 $\nu_{x,t}$ 解释为\eqref{align:ch17-1731}的一种新型弱解：

\begin{definition}[测度值解]\label{definition:ch17-测度值解}
守恒律方程组\eqref{align:ch17-1731}在 $\mathbb{R}^2$ 的开子集 $\Omega$ 上的测度值解，是一可测的概率测度族 $\{\nu_{x,t}:(x,t)\in\Omega\}$，它在 $\Omega$ 上的分布意义下满足\eqref{align:ch17-1734}。
\end{definition}

显然，\eqref{align:ch17-1731}的任意传统弱解 $U\in L^\infty(\Omega;\mathbb{R}^n)$ 都可与测度值解 $\nu_{x,t}=\delta_{U(x,t)}$ 等同。然而，测度值解类无疑比传统解类宽泛得多。例如，若 $U$ 与 $V$ 是\eqref{align:ch17-1731}在 $L^\infty(\Omega;\mathbb{R}^n)$ 中的任意两个传统解，则对任意固定的 $\alpha\in(0,1)$，
\begin{align}
\nu_{x,t}=\alpha\delta_{U(x,t)}+(1-\alpha)\delta_{V(x,t)}
\label{align:ch17-1735}
\end{align}
定义了一个非传统的测度值解。

乍看之下，测度值解这一概念可能宽泛到不具现实意义。然而，放弃“解应在每一点 $(x,t)$ 对状态向量赋以确定值”这一前提，恰好提供了描述一类物理现象的手段，例如相变：在宏观层面，相的混合物可占据时空同一点。我们在此不展开这些想法，而是把测度值解视作构造传统解的垫脚石。

可容性概念自然地从传统解推广到测度值解。若相对于\eqref{align:ch17-1731}的熵--熵流对 $(\eta,q)$，$\Omega$ 上的测度值解 $\nu_{x,t}$ 满足
\begin{align}
\partial_t\langle\nu_{x,t},\eta(U)\rangle+\partial_x\langle\nu_{x,t},q(U)\rangle\leqslant 0,
\label{align:ch17-1736}
\end{align}
在 $\Omega$ 上的分布意义下，则称它满足熵可容性条件。

回到前面的例子，设 $\nu_{x,t}$ 由\eqref{align:ch17-1733}的解序列 $\{U_{\mu_j}\}$ 生成。若 $(\eta,q)$ 是\eqref{align:ch17-1731}的任意熵--熵流对，将\eqref{align:ch17-1733}乘以 $D\eta(U_\mu)$ 并利用\eqref{align:ch07-熵流对定义}，
可得
\begin{align}
\partial_t\eta(U_\mu)+\partial_x q(U_\mu)=\mu\,\partial_x^2\eta(U_\mu)-\mu\,\partial_xU_\mu^{\top}\mathrm{D}^2\eta(U_\mu)\partial_xU_\mu.
\label{align:ch17-1737}
\end{align}
特别地，当 $\eta$ 为凸时，\eqref{align:ch17-1737}右端最后一项非正。于是我们断言：由相对于\eqref{align:ch17-1733}的消失黏性方法构造的、\eqref{align:ch17-1731}的任意测度值解 $\nu_{x,t}$，对任意 $\eta$ 为凸的熵--熵流对 $(\eta,q)$，都满足熵可容性条件\eqref{align:ch17-1736}。

不应认为可容性本身足以把测度值解约化为传统解。应当指出：当两个传统解 $U$ 与 $V$ 对某熵--熵流对 $(\eta,q)$ 满足熵可容性条件时，由\eqref{align:ch17-1735}定义的非传统测度值解 $\nu_{x,t}$ 同样满足该条件。另一方面，在测度值解的框架下，可容性也可能成为唯一性与稳定性的动因。沿此方向，已有结论（文献见\cref{section:ch17-notes}）表明：在上半平面，任一标量守恒律的测度值解 $\nu_{x,t}$，若对所有凸熵--熵流对满足熵可容性条件，且初值为 Dirac 质量 $\nu_{x,0}=\delta_{u_0(x)}$（其中 $u_0\in L^\infty(-\infty,\infty)$），则必约化为传统解，即 $\nu_{x,t}=\delta_{u(x,t)}$，其中 $u$ 是以 $u(x,0)=u_0(x)$ 为初值的该守恒律唯一可容解。特别地，这蕴含：对标量守恒律而言，由消失黏性方法以传统初值构造的任意测度值解都约化为传统解。

回到方程组\eqref{align:ch17-1731}，下面将勾勒一个纲领，用以验证：与某近似解序列 $\{U_j\}$ 相联系的 Young 测度族 $\{\nu_{x,t}:(x,t)\in\Omega\}$ 所诱导的测度值解约化为传统解。该纲领随后将针对特殊方程组实现。如\cref{section:ch01-快速振荡}已指出的，
当\eqref{align:ch17-1731}为双曲型时，近似解可能发展出持续的快速振荡，从而妨碍序列 $\{U_j\}$ 的强收敛。因此，除非近似解以某种方式体现出淬灭振荡的机制，否则我们的事业注定失败。从补偿紧致性理论的立场看，这种机制体现为条件
\begin{align}
\partial_t\eta(U_j)+\partial_x q(U_j)\subset\text{紧集于 }W^{-1,2}_{\mathrm{loc}}(\Omega),
\label{align:ch17-1738}
\end{align}
它对\eqref{align:ch17-1731}的任意熵--熵流对 $(\eta,q)$ 成立。

为看清\eqref{align:ch17-1738}的含义，任取两个熵--熵流对 $(\eta_1,q_1)$ 与 $(\eta_2,q_2)$。当 $j\to\infty$ 时，序列 $\{\eta_1(U_j)\}$、$\{\eta_2(U_j)\}$、$\{q_1(U_j)\}$、$\{q_2(U_j)\}$ 分别收敛于 $\bar\eta_1=\langle\nu,\eta_1\rangle$、$\bar\eta_2=\langle\nu,\eta_2\rangle$、$\bar q_1=\langle\nu,q_1\rangle$、$\bar q_2=\langle\nu,q_2\rangle$，其中为简洁记 $\nu_{x,t}=\nu$。由\eqref{align:ch17-1738}，$\operatorname{div}(q_2(U_j),\eta_2(U_j))$ 与 $\operatorname{curl}(\eta_1(U_j),-q_1(U_j))$ 均含于 $W^{-1,2}_{\mathrm{loc}}(\Omega)$ 的紧集。于是由\cref{theorem:ch17-divcurl引理}，
\begin{align}
\eta_1(U_j)q_2(U_j)-\eta_2(U_j)q_1(U_j)\rightharpoonup\bar\eta_1\bar q_2-\bar\eta_2\bar q_1,\qquad j\to\infty,
\label{align:ch17-1739}
\end{align}
在 $L^\infty(\Omega)$ 弱$^*$ 意义下，等价地
\begin{align}
\langle\nu,\eta_1\rangle\langle\nu,q_2\rangle-\langle\nu,\eta_2\rangle\langle\nu,q_1\rangle=\langle\nu,\eta_1q_2-\eta_2q_1\rangle.
\label{align:ch17-17310}
\end{align}

我们的计划是：对策略性选取的熵--熵流对使用\eqref{align:ch17-17310}，以证明 Young 测度 $\nu$ 的支撑被限制在单点上。显然，只有当具备按预定规格构造熵--熵流对的灵活性时，这一纲领才有成功的可能。在一切实际用途上，这一要求把该方法的适用范围限制于标量守恒律、两个守恒律的方程组，以及那些配备了丰富熵族（见\cref{section:ch07-熵熵流对}）的、多于两个守恒律的特殊方程组。
另一方面，该方法对构造格式具有相当的灵活性，因为它只要求近似解满足\eqref{align:ch17-1738}。

作为说明，我们在“方程组\eqref{align:ch17-1731}配备一致凸熵、$\Omega$ 为上半平面、近似解序列 $\{U_j\}$ 由消失黏性方法生成（$U_j=U_{\mu_j}$，其中 $U_\mu$ 是\eqref{align:ch17-1733}在上半平面上以
\begin{align}
U(x,0)=U_{0\mu}(x),\qquad-\infty<x<\infty,
\label{align:ch17-17311}
\end{align}
为初值、且 $U_{0\mu}$ 含于 $L^\infty(-\infty,\infty)\cap L^2(-\infty,\infty)$ 的一个有界集）”这一假设下验证\eqref{align:ch17-1738}。

设 $\eta$ 是一致凸熵，因而 $\mathrm{D}^2\eta(U)$ 正定。我们可设 $0\leqslant\eta(U)\leqslant c|U|^2$，因为否则只需用熵 $\eta^*(U)=\eta(U)-\eta(0)-\mathrm{D}\eta(0)U$ 替换 $\eta$。将\eqref{align:ch17-1737}在上半平面积分，得到估计
\begin{align}
\mu\int_0^\infty\int_{-\infty}^\infty|\partial_xU_\mu(x,t)|^2\,dx\,dt\leqslant a,
\label{align:ch17-17312}
\end{align}
其中 $a$ 与 $\mu$ 无关。

现任取一个不必为凸的熵--熵流对 $(\eta,q)$，并在上半平面中取定某开有界子集 $\Omega$。考察\eqref{align:ch17-1737}：左端在 $W^{-1,p}(\Omega)$（任意 $1\leqslant p<\infty$）中有界；右端是两项之和——由\eqref{align:ch17-17312}，第一项当 $\mu\downarrow0$ 时在 $W^{-1,2}(\Omega)$ 中趋于零，因而含于 $W^{-1,2}(\Omega)$ 的一个紧集；第二项含于 $M(\Omega)$ 的一个有界集，这同样由\eqref{align:ch17-17312}得来。因此，\eqref{align:ch17-1738}由\cref{lemma:ch17-插值紧致引理}推出。

\end{document}
